\section{Discrete groups}\label{sec:discrete}

For locally compact groups that are not compact, and in particular for
infinite discrete groups, the universal minimal flow is a large,
non-metrizable space. For countable discrete groups its topology is
nevertheless determined up to one cardinal parameter.

\begin{remark}[Size]\label{rem:umf-size}
\uses{thm:veech}
If $G$ is locally compact and not compact, then $M(G)$ is not metrizable
\cite{GlasnerWeiss2003,GiordanoPestov2007}; compare
Remark~\ref{rem:lc-metrizable}. For infinite discrete groups this is due to
Ellis, and Kechris, Pestov and Todorcevic extended it to locally compact
groups that are not compact \cite{GlasnerGutman2012}. By Theorem~\ref{thm:veech}, $G$ acts freely
on it.
\end{remark}

\begin{example}\label{ex:discrete-ambit}
\uses{def:samuel}
If $G$ is discrete, every bounded function on $G$ is right uniformly
continuous, so $\mathrm{RUC}_b(G)=\ell^\infty(G)$, and $S(G)$ is the
Stone--Čech compactification $\beta G$.
\end{example}

\begin{proposition}\label{prop:extremely-disconnected}
\uses{ex:discrete-ambit, prop:minimal-ideals}
If $G$ is discrete, then $M(G)$ is extremely disconnected: the closure of
every open subset is open.
\end{proposition}

\begin{proof}
The space $\beta G$ is extremely disconnected, and by
Proposition~\ref{prop:minimal-ideals} there is a retraction
$r\colon\beta G\to M$ onto a minimal subflow $M\cong M(G)$. Let $U\subseteq M$
be open and $V=r^{-1}(U)$, so that $U\subseteq V$ and $r(V)=U$. Then
$\overline{U}\subseteq\overline{V}\cap M\subseteq r(\overline{V})
\subseteq\overline{r(V)}=\overline{U}$, where the closure of $U$ is taken in
$M$. Hence $\overline{U}=\overline{V}\cap M$, which is open in $M$ because
$\overline{V}$ is open in $\beta G$.
\end{proof}

\begin{theorem}[Balcar--Błaszczyk \cite{BalcarBlaszczyk1990}]\label{thm:balcar-blaszczyk}
Let $B$ be an infinite Boolean algebra and $H$ a countable group of
automorphisms of $B$ such that $(B,H)$ is \emph{minimal}: for every nonzero
$u\in B$ there are $h_1,\dots,h_n\in H$ with $h_1(u)\vee\dots\vee h_n(u)=1$.
Then the completion of $B$ is isomorphic to the completion of the free Boolean
algebra on $\pi(B)$ generators, where $\pi(B)$ is the least cardinality of an
order-dense subset of $B\setminus\{0\}$.
\end{theorem}

\begin{remark}\label{rem:balcar-blaszczyk}
\uses{thm:balcar-blaszczyk}
Minimality of $(B,H)$ means that the dual action of $H$ on the Stone space of
$B$ is minimal. The countability of $H$ cannot be dropped: a homogeneous
Boolean algebra with all its automorphisms is a minimal system
\cite{BalcarBlaszczyk1990}.
\end{remark}

\begin{corollary}\label{cor:discrete-cohen}
\uses{thm:balcar-blaszczyk, prop:extremely-disconnected, thm:veech}
Let $G$ be a countably infinite discrete group. Then $M(G)$ is homeomorphic to
the Stone space of the completion of a free Boolean algebra; equivalently, to
the Gleason cover (absolute) of a Cantor cube $\{0,1\}^\kappa$.
\end{corollary}

\begin{proof}
By Proposition~\ref{prop:extremely-disconnected}, the clopen subsets of $M(G)$
form a complete Boolean algebra $B$, whose Stone space is $M(G)$. It is
infinite, because $G$ is infinite and acts freely on $M(G)$
(Theorem~\ref{thm:veech}). The group $G$ acts on $B$ by automorphisms, and
minimality of $M(G)$ says that for every nonzero $b\in B$ finitely many
translates of $b$ cover $M(G)$, which is minimality of $(B,G)$. Since $B$ is
complete, Theorem~\ref{thm:balcar-blaszczyk} applies to $B$ itself. The
completion of the free Boolean algebra on $\kappa$ generators is the algebra
of regular open subsets of $\{0,1\}^\kappa$, whose Stone space is the Gleason
cover of $\{0,1\}^\kappa$.
\end{proof}

\begin{example}[The integers]\label{ex:umf-integers}
\uses{cor:discrete-cohen, prop:minimal-ideals}
For $G=\mathbb{Z}$, the cardinal is the continuum. Balcar and Błaszczyk show
that every minimal subflow of $\beta\mathbb{Z}$, with the action extending
the shift, is homeomorphic to the Gleason cover of $\{0,1\}^{\mathfrak{c}}$
\cite{BalcarBlaszczyk1990}. By Proposition~\ref{prop:minimal-ideals}, so is
$M(\mathbb{Z})$: it is the Stone space of the Boolean algebra of regular open
subsets of $\{0,1\}^{\mathfrak{c}}$; see also \cite{BartosovaKwiatkowska2019}.
\end{example}
